\documentclass{article}
\PassOptionsToPackage{round,authoryear}{natbib}
\usepackage[dblblindworkshop]{neurips_2026}
\workshoptitle{3rd Workshop on Attributing Model Behavior at Scale: Data Attribution and Provenance (ATTRIB 2026)}

\usepackage[utf8]{inputenc}
\usepackage[T1]{fontenc}
\usepackage{hyperref}
\usepackage{url}
\usepackage{booktabs}
\usepackage{amsmath,amssymb}
\usepackage{amsfonts}
\usepackage{microtype}
\usepackage{graphicx}

\newtheorem{lemma}{Lemma}

\title{Correlated Failure as an Indirect Signature of Shared Model \\
Provenance: A Population-Scale Complement to Instance-Level Attribution}

\author{
  Anonymous Author(s) \\
  Affiliation \\
  Address \\
  \texttt{email}
}

\begin{document}

\maketitle

\begin{abstract}
Contributive and corroborative attribution methods answer ``which training data explains this
output,'' one output at a time, and they need access to training data or model internals that is
routinely unavailable for open leaderboard submissions. We describe a training-data-free,
population-scale signal that is available even when neither is: after conditioning on the exact,
fit-free null implied by Rasch sufficiency, $1{,}228$--$1{,}362$ open models per benchmark share
residual failure correlation $2.8$--$11.8\times$ their null, concentrated in a handful of latent
dimensions rather than one shared factor, and this signal survives removing exact near-duplicate
submissions. We could not decompose this signal against declared model lineage --- a
pre-registered check found only $23.3\%$ of models had a resolvable \texttt{base\_model} field, so
we dropped the analysis rather than estimate it on a biased subset. We think that failure is itself
informative: it is a concrete instance of exactly the provenance-opacity problem this workshop's
community is positioned to attack, and we pose the residual correlation matrix as an idea-track
candidate for what a black-box, aggregate provenance signal could look like when nothing about
training data is disclosed.
\end{abstract}

\section{The gap we are pointing at}

Contributive attribution estimates the causal effect of specific training examples on a specific
generation; corroborative attribution finds sources that textually support a specific output. Both
operate at the level of one model, one output. Neither is designed for --- and as far as we know,
neither has been tried on --- the regime of an open leaderboard with thousands of independently
submitted models, most of unknown or undisclosed lineage, where the question is not ``which
example explains this generation'' but ``which models share unstated ancestry with which other
models, and can that be seen at all without the training data.''

We ran into this question directly, from the measurement side rather than the attribution side,
while trying to characterize correlated failure between open language models on standard
multiple-choice benchmarks. We report what we found because we think the gap between what we
needed and what currently exists is itself a useful data point for this workshop, even though we
do not close it.

\section{A population-scale, training-data-free residual signal}

Let $F_{im}=1$ iff model $i$ fails item $m$. The row and column margins of $F$ are the jointly
sufficient statistics of the Rasch model, so conditioning on both --- sampled here with curveball
\citep{strona2014}, verified against enumeration on small fibres --- gives an exact, fit-free null
under which mean-level co-failure is provably degenerate: it is a deterministic function of item
difficulty alone, independent of any shared behavior between models
\citep{schluter1984,gotelli2000,cronbach1951}. At that null, on the archived Open LLM Leaderboard
(five benchmarks, $1{,}228$--$1{,}362$ models each, harvested from result files at zero inference
cost), the residual correlation structure that remains after removing the difficulty degeneracy is
$2.8$--$11.8\times$ its null value, and spectrally concentrated: on ARC-Challenge, six eigenvalues
exceed the null's spectral edge and the leading eigenvalue explains only $54\%$ of squared residual
mass, the signature of several latent dimensions rather than one shared factor. This is not a
duplicate-submission artifact: clustering models at agreement $1.0000$ and keeping one
representative per cluster removes $900$ of $1{,}362$ ARC models and moves the calibrated
participation-ratio ratio from $0.052$ only to $0.058$. Whatever produces this structure survives
the crudest form of provenance filtering available to us.

That last clause is doing real work: it is the crudest form available, because we do not have
access to a better one. We do not observe pretraining corpora, fine-tuning data, or distillation
relationships for the large majority of these models. A leaderboard submission carries an
optional, self-reported \texttt{base\_model} field, and in our sample it was resolvable for only
$23.3\%$ of models --- below a threshold we set before looking at the data specifically to decide
whether a lineage decomposition was trustworthy. It fell short, so we dropped that analysis rather
than report it on a biased subset. We are not claiming the residual correlation \emph{is} shared
provenance; the honest statement is narrower and, we think, more useful to this audience: standard
attribution machinery cannot currently tell us either way, because it needs exactly the access this
population does not provide.

\section{Why this might be attribution's problem, not measurement's}

The residual structure has three properties that make it look like a provenance signal rather than
a pure behavioral coincidence. First, it is not explained by discrimination heterogeneity alone: a
two-parameter item-response simulation with realistic discrimination variance produces residual
correlation of at most $0.076$ (rms), against an observed $0.248$ on ARC. Second, it is not
explained by literal duplicate submissions, as above. Third, its dimensionality is stable and
benchmark-dependent ($4$--$6$ latent dimensions on four benchmarks, $12$ on HellaSwag), which is
the kind of structure one would expect from a handful of dominant pretraining lineages or shared
distillation teachers producing correlated errors along a few behavioral axes, rather than from
either a single dominant lineage or from fully independent training runs.

None of this is proof of shared provenance --- discrimination heterogeneity, unmodeled ability
dimensions, and genuinely convergent training choices are alternative explanations we cannot rule
out with this design, and our own limitations section says so plainly. What we can say is that the
exact conditional null gives a population-scale test that is exactly invariant to item difficulty
by construction (Lemma~\ref{lem:mean}), costs no training-data access, and is sensitive to
structure at a scale (thousands of models) where per-example causal attribution is not currently
run at all.

\begin{lemma}[Mean co-failure depends only on item margins]
\label{lem:mean}
$O = \frac{1}{MN(N-1)}\sum_m(c_m^2-c_m)$, where $c_m$ is the number of models failing item $m$;
hence any randomization preserving item margins leaves $O$ exactly invariant.
\end{lemma}

\section{An idea-track question}

We pose this as a question rather than a method: can the residual eigenvectors of a
margin-conditioned co-failure matrix be matched --- even loosely, even just as a screening tool ---
against known facts about model families, base-model announcements, or public distillation claims,
to turn an aggregate statistical signal into something closer to an attribution claim? If so, this
would be a training-data-free, black-box complement to contributive and corroborative methods,
useful specifically in the regime where those methods' preconditions (access to training data,
identified individual outputs to explain) do not hold, such as regulatory or journalistic audits of
a leaderboard's actual diversity of origin. If not --- if the residual structure is better
explained by convergent unimodal training choices than by shared lineage, which our own spectral
diagnostics do not rule out --- that would itself be a useful negative result about how much
population-scale behavioral correlation the field should be willing to read as a provenance signal
at all. We do not know which is true, and we think that is the right question for this venue rather
than for a measurement paper on its own.

\section{Limitations}

We did not attempt any form of causal or corroborative attribution ourselves; this paper reports a
measurement and a gap, not an attribution method. Discrimination heterogeneity and unmodeled ability
dimensions are alternative explanations for the residual structure that our design cannot rule out.
Claims are restricted to the evaluated open-model ecosystem on multiple-choice benchmarks. Code,
data, and the pre-registration are archived and linked in the full paper.

\begin{thebibliography}{9}
\small
\bibitem[Cronbach(1951)]{cronbach1951} L.~J. Cronbach. Coefficient alpha and the internal structure of tests. \emph{Psychometrika}, 16:297--334, 1951.
\bibitem[Gotelli(2000)]{gotelli2000} N.~J. Gotelli. Null model analysis of species co-occurrence patterns. \emph{Ecology}, 81(9):2606--2621, 2000.
\bibitem[Schluter(1984)]{schluter1984} D.~Schluter. A variance test for detecting species associations, with some example applications. \emph{Ecology}, 65:998--1005, 1984.
\bibitem[Strona et al.(2014)]{strona2014} G.~Strona, D.~Nappo, F.~Boccacci, S.~Fattorini, and J.~San-Miguel-Ayanz. A fast and unbiased procedure to randomize ecological binary matrices with fixed row and column totals. \emph{Nature Communications}, 5:4114, 2014.
\end{thebibliography}

\end{document}
